% Generated draft: requires mathematical and source review.
\documentclass[11pt]{article}
\usepackage[margin=25mm]{geometry}
\usepackage{fontspec}
\setmainfont{Latin Modern Roman}
\usepackage{amsmath,amssymb,mathtools,unicode-math}
\setmathfont{Latin Modern Math}
\usepackage{enumitem,needspace,xurl,hyperref}
\setlength{\parindent}{0pt}
\setlength{\parskip}{5pt}
\setlength{\emergencystretch}{3em}
\newcommand{\N}{\mathbb{N}}
\newcommand{\Z}{\mathbb{Z}}
\newcommand{\Q}{\mathbb{Q}}
\newcommand{\R}{\mathbb{R}}
\newcommand{\C}{\mathbb{C}}
\newcommand{\F}{\mathbb{F}}
\newcommand{\A}{\mathbb{A}}
\newcommand{\PP}{\mathbb{P}}
\newcommand{\E}{\mathbb{E}}
\newcommand{\eps}{\varepsilon}
\newcommand{\dd}{\,\mathrm{d}}
\begin{document}
\section*{UnsolvedMath ID: 8700030}
\textbf{Source problem code:} catalogue entry
\subsection{Definitions and mathematical statement}
Fix a prime \(p\), work in \(\C_p\) with \(|p|_p=p^{-1}\), and use a fixed embedding of algebraic numbers. The power-series conventions are
\[
\log_p(1+u)=\sum_{k\ge1}\frac{(-1)^{k+1}u^k}{k}\quad(|u|_p<1),
\qquad
\exp_p(z)=\sum_{k\ge0}\frac{z^k}{k!}
\quad(|z|_p<p^{-1/(p-1)}).
\]
Let \(\alpha\ne0\) be algebraic in the logarithm domain and let \(\beta\) be cubic algebraic in the intended p-adic field. A conventional degree reading is \([\Q(\beta):\Q]=3\); whether \(\beta\) must lie in \(\Q_p\) rather than a larger p-adic field is not explicit.

Put \(L=\log_p\alpha\). When both exponential arguments are in their domains, define
\[
E_1=\alpha^\beta=\exp_p(\beta L),\qquad
E_2=\alpha^{\beta^2}=\exp_p(\beta^2L).
\]
The intended algebraic-independence assertion is
\[
\forall P\in\Q[U,V]\setminus\{0\},\qquad P(E_1,E_2)\ne0.
\]
The source assumes only convergence for \(\beta L\). This does not by itself guarantee convergence for \(\beta^2L\) when \(|\beta|_p>1\); both arguments need a justified domain condition. A hypothesis such as \(|\beta|_p\le1\) would suffice for that step, but is not stated and is not silently inserted.

There is also an explicit degeneracy: \(\alpha=1\) satisfies the written hypotheses and gives \(E_1=E_2=1\), contradicting algebraic independence. More generally \(L=0\) gives this obstruction. A meaningful intended version therefore requires \(\log_p\alpha\ne0\), as well as definedness of both exponentials, or an equivalent source qualification. These are flagged corrections, not replacements of the supplied literal statement. The research summary calls the intended p-adic cubic independence problem open and compares it with a classical complex Gel'fond result, without resolving these missing hypotheses.
\subsection{Short English statement}
The intended conjecture asserts algebraic independence of two p-adic exponentials with cubic exponent parameter. As written it allows \(\alpha=1\) and does not guarantee that the second exponential converges, so its nondegeneracy and domain hypotheses need repair.
\subsection{Sources}
\begin{itemize}
\item Local UnsolvedMath registry, permanent ID \(8700030\), exact source code \texttt{catalogue entry}, ``Conjecture 3.12 --- p-adic analog of an algebraic independence result of Gel'fond.'' Proposer Michel Waldschmidt.
\item Waldschmidt, \emph{Open Diophantine Problems} (2004), Conjecture 3.12; \url{https://arxiv.org/abs/math/0312440}. Source-TeX extraction; recorded access 2026-07-29.
\item Supplied OPEN-TRIAGE research summary describes the intended p-adic question as unresolved and mentions the classical complex Gel'fond theorem, but gives no further precise citation or missing-domain clarification.
\end{itemize}
\end{document}
